% Pursuing Stacks (cotes 134-2 to 134-8): archive cote 134-2, site page 14 (PDF page 15).
% letter to Quillen. Draft compared: ps1.tex:3-65.
% AI-transcribed by Claude Opus 5.5 (claude-opus-5-5, adaptive thinking, effort high), Message Batches API, batch ps1 (msgbatch_01Rv5hCaNTtUt4JHx1yqoHPL), 2026-11-12; not yet checked by eye.
\scanpage{134-2/14}
\mspage{\hand{1}}

Chapter I : T A K E \unsure{-} O F F

A letter to ....

\hand{(}followed by an \over{\unsure{a}}{A}ppendix: Three letters to Larry Breen\hand{)}

% NOTE: from here on the text stands on a strip of paper pasted onto the page; the upper edge of the strip is visible
\margin{1)}

\margin{Les Aumettes 19.2.1983}

\margin{\unsure{\#}}

Dear \deleted{\unsure{Daniel},} ---

\hand{\S} 1. Last year Ronnie Brown from Bangor sent me a heap of reprints and pre-%
prints by him and a group of friends\hand{,} on various foundational matters of
homotopical algebra. I did not really dig through any of this, as I kind of
lost contact with the technicalities of this kind (I was never too familiar
with the homotopy techniques\inserted{anyhow,}\hand{/} I confess) - but this reminded me of a few
letters I had exchanged with Larry Breen in 1975\hand{(1)}, where I had developped
an outline of a \inserted{program for a} kind of ``topological algebra'', viewed as a synthesis of
\unsure{homotopical} and homological algebra, with special emphasis on topoi -
most of the basic intuitions in this program arising from various backgrou\hand{nd}
in algebraic geometry. \hand{\unsure{(}}Some of those intuitions we discussed\hand{,} I believe\hand{,}
at IHES eight or nine years before\hand{,} \inserted{then, \unsure{at a time} when\dots?} \deleted{at a time when you had just written
up your nice ideas on axiomatic homotopical algebra, published since in
Spinger's Lecture Notes.}\hand{/} I write you under the assumption that you have not
entirely lost interest for those foundational questions you were looking at
more than fifteen years ago. One thing which strikes me, is that (as far
as I know) there has not been any substantial progress since - it looks
to me that an\xout{\illegible{1}} understanding of the basic structures underlying homoto-%
py theory, or even homological algebra only, is still lacking - probably
because the few people who have a wide enough background and perspective
enabling them to feel the main questions, are devoting their energies to
things which seem more directly rewarding. Maybe \xout{\unsure{there}} even a wind of dis-%
repute for any foundational matters whatever is blowing nowadays \xout{\unsure{!}}\hand{(2)}\xout{\unsure{In}}\hand{! In}
this reppect, what seems to me even more striking than the lack of proper
foundations \over{of}{\unsure{for}} homological and homotopical algebra, is the absence I dare-%
say of proper foundations for topology itself ! I am thinking here mainly
of the development of a context of ``tame'' topology, which (I am convinced\hand{)}
would have on the everyday techhique of geometric topology (\xout{\unsure{in}} I use this
expression in contrast to the topology of use for analysts) a comparable
impact or even a greater one, than the introduction of the point of view of
schemes had on algebraic geometry\hand{(3)}.The \inserted[typed]{psychological} drawback here I believe is not any-%
thing like messyness, as for homological and homotopical algrebra\inserted{(or for schemes)}\hand{,} but mer-%
\margin{-}ly the inrooted intertia which \xout{\unsure{goes against}} prevents \xout{\illegible{1}} us so stub\hand{born}-%
ly from looking innocently, with fresh eyes, upon things, without being
dulled and emprisoned by standing habits of thought, going with a familiar
context - \uline{too} familiar a context ! The task of working out the foundati-%
ons of tame topology\hand{,} and a \xout{\unsure{theory}} corresponding structure theory for ``stra-%
tified (tame) spaces\hand{\unsure{(4)}}\hand{,} seems to me a lot more urgent and exciting still \inserted{than} \xout{\unsure{th}}

\margin{14}
